\SourceLeaf{335}%
\Article{34}
Of the seven coefficients (1), (2), and so on, the last is constant, with
\[
\log(7)=7{,}6287228032(-20)
\]
and
\[
\log3(7)=8{,}1058440580(-20)
\]
Once specific values are chosen for the dimensions of the ellipsoid, the other coefficients become functions of the latitude \(B\). They can therefore be tabulated with \(B\) as the argument. With such a table at hand, computation by this method is as convenient for the ellipsoid as it is for the sphere.

At the end of this memoir I include such a table for the zone from \(51^0\) to \(54^0\), in which the values of \(B\) advance by one minute at a time. I make the following remarks about it.

The table depends on the parameters of the ellipsoid only insofar as it uses a particular flattening, or a particular value of \(e\). This is the value obtained in the latest calculation by \textsc{Bessel}, and it is also the basis of the table appended to the first memoir (see Art. 5). To make the numerical value of \(a\) depend only on the flattening, the unit chosen is the ten-millionth part of the Earth’s meridian, in place of the toise or any other arbitrary measure. Thus \(a\) can be found directly from \(e\) by means of the equation
\[
\begin{aligned}
\pi a(1-\tfrac14ee-\frac{1\mathord{.}3}{4\mathord{.}16}e^4-\frac{1\mathord{.}3\mathord{.}15}{4\mathord{.}16\mathord{.}36}e^6\\
{}-\frac{1\mathord{.}3\mathord{.}15\mathord{.}35}{4\mathord{.}16\mathord{.}36\mathord{.}64}e^8
-\frac{1\mathord{.}3\mathord{.}15\mathord{.}35\mathord{.}63}{4\mathord{.}16\mathord{.}36\mathord{.}64\mathord{.}100}e^{10}-\text{and so on})\\
=20000000
\end{aligned}
\]
whose pattern is evident, or by means of the equivalent equation
\[
a=\frac{20000000}{\pi}(1+\tfrac14ee+\tfrac7{64}e^4+\tfrac{15}{256}e^6+\tfrac{579}{16384}\mathord{.}e^8+\tfrac{1515}{65536}e^{10}+\text{and so on})
\]
This gives the following results for that value of \(e\).
\[
\begin{aligned}
a&=6376851{,}447\\
\log a&=6{,}8046062999
\end{aligned}
\]
Of course, to use our table one must first express \(r\) in the same unit. To make this conversion, one must, using the value given by
